% !TEX root = ../../main.tex
% C3.4 - Frame sampling (translated). Matches workers/sampling.py (i mod N = 0; baseline N=10, throughput N=20 default).
\section{Frame Sampling}
\label{sec:3.4}
\label{sec:3.4.3}

The project's videos have about 60 frames per second, so consecutive frames are about 17 milliseconds apart and almost identical, and running the models on every frame is costly for little new information. Frames are sampled uniformly by index: frame $i$ is selected with sampling interval $N$ if and only if
\begin{equation}
    i \bmod N = 0,
    \label{eq:sampling}
\end{equation}
and it keeps its original index and timestamp, so that the time at which a person appears follows the camera's own timeline. With a source of $f$ frames per second, two processed frames are $N/f$ seconds apart: about 0.17 seconds with $N = 10$ and 0.33 seconds with $N = 20$.

The interval is a trade-off. Sampling precedes the detector, so model runs drop about $N$ times, but every frame must still be decoded, since most H.264 frames depend on earlier ones (Section~\ref{sec:3.7.1}). A larger $N$ misses brief appearances more easily, since a track needs two observations, and lets people move further between observations, making IoU association harder. The system provides two presets: \emph{baseline} with $N = 10$, which favours track quality, and \emph{throughput} with $N = 20$, which favours speed and was chosen as the default based on measurements on the target machine (Section~\ref{sec:8.5}). Each job records the value of $N$ it used so that results can be compared and reproduced.
